\subsection{History and scope}\label{hist:history}

Penrose's 1973 argument connected gravitational collapse and cosmic censorship
with a lower bound for the mass in terms of black-hole area
\cite{penrose1973}.  In the asymptotically flat, time-symmetric setting,
Huisken and Ilmanen proved the three-dimensional inequality for the area of
each connected outermost minimal component by weak inverse mean curvature
flow \cite{huiskenilmanen2001}.  Bray's conformal flow proved the bound for
the total area of a possibly disconnected outermost minimal boundary
\cite{bray2001}; Bray and Lee extended the Riemannian inequality to dimensions
less than eight \cite{braylee2009}.  These are scalar-curvature theorems for
Riemannian metrics.  Recovering an arbitrary original second fundamental
form in an equality statement requires further arguments.

For cosmological constant $-3$, the Schwarzschild--anti-de Sitter horizon
relation is $M=(r_h+r_h^3)/2$.  It leads to the asymptotically hyperbolic
Penrose bound, whose mass convention comes from the metric mass integrals
of Wang and Chru\'sciel--Herzlich \cite{wang2001,chruscielherzlich2003}.
Earlier Riemannian results include the graphical estimates of Dahl,
Gicquaud and Sakovich \cite{dahlgicquaudsakovich2013}, the sharp Inequality
of de Lima and Gir\~ao for balanced hyperbolic graphs with scalar curvature
at least $-6$ and a minimal horizon in a horizontal totally geodesic slice
met orthogonally, whose projection is star-shaped and mean convex
\cite[Theorem~1.2]{delimagirao2016},
and Ambrozio's Theorem for sufficiently small perturbations of a fixed
positive-mass Schwarzschild--anti-de Sitter metric, with minimal boundary
and scalar curvature at least $-6$ \cite{ambrozio2015}.
Lee and Neves treat a different asymptotically locally hyperbolic regime,
using a nonpositive supremum of the mass aspect and a boundary component
with the genus of conformal infinity \cite{leeneves2015}.
For static systems, Borghini, Fogagnolo and Pinamonti characterize
finite domains attaining equality in the substatic Heintze--Karcher
inequality as warped products
\cite[Theorem~1.1]{borghinifogagnolopinamonti2024}.
We use this finite-domain result after proving that the equality data
force the static deficit to vanish, and then prove the global extension.
The flow issue at infinity is substantive: Neves constructed an
asymptotically hyperbolic inverse mean curvature flow lacking the convergence
needed for the direct Hawking-mass argument \cite{neves2010}.

There are also perturbative results involving the second fundamental form.
Khuri and Kopi\'nski prove the inequality for small maximal vacuum conformal
perturbations of Schwarzschild--anti-de Sitter data under a quantitative
comparison between a weighted bulk norm and a boundary-related norm of the
seed tensor \cite{khurikopinski2023}.
The earlier method manuscript \cite{internalads} treats a fixed seed and a
fixed local maximal vacuum conformal branch; its quadratic-null directions
require a fourth-order calculation.  The asymptotically flat manuscript
\cite{internalflat} develops coupled deformation and original-data rigidity
methods at zero cosmological constant.  We use these as methodological
antecedents and supply the arguments needed at the present hypotheses.

\subsection{Significance and technical contributions}\label{hist:significance}

We establish the three-dimensional nonmaximal asymptotically anti-de Sitter
spacetime Penrose Inequality for the precise class specified below.
Its numerical assertion uses the full enclosing-area infimum
of a possibly disconnected weakly future outer trapped
boundary and the Lorentz norm of the four \emph{metric} fluxes.
In the connected, outermost, outer area-minimizing future marginally
outer trapped subclass, its equality assertion recovers the original pair
$(g,K)$ as a smooth proper spacelike hypersurface in the smooth maximal
extension of Schwarzschild--anti-de Sitter spacetime of parameter exactly
$M=m_{\mathrm{AH}}$, with a full future horizon section (the bifurcation
sphere is allowed) and an end reaching conformal infinity.
This resolves the asymptotically anti-de Sitter spacetime Penrose Conjecture
in this initial-data class, including its Riemannian specialization and
original-data rigidity on the stated connected horizon subclass.
The mass match concerns the designated initial
slice; spatial chart covariance supplies no invariance under arbitrary
changes of that slice.  The assertions are initial-data statements and
make no claim to prove cosmic censorship or an evolution theorem.

Three parts of the construction have uses beyond the model calculation.
First, the coupled deformation combines weak trapping barriers with
estimates that distinguish the constants required uniformly in the large
profile parameter from those needed only after that parameter is fixed.
A scalar regularity argument for a measurable rank-one axial coefficient
and a bounded-input solution operator provide the analytic ingredients
for the degree construction.  Second, a primitive correcting the spatially
varying lapse cancels mixed derivatives in the hyperbolic deformation;
the resulting weighted divergence estimate controls all four mass fluxes
and retains the full enclosing-area comparison.  Third, equality is tested
on the original constraints.  The first stationarity argument determines
the horizon surface gravity, and a causal adjoint with analysis across its
null set produces the complete static quotient.  Weighted improvement of
its end evaluates the static Heintze--Karcher deficit in terms of the
original mass. The equality data force its nonnegative interior defect
to vanish. A finite horizon collar then attains exact equality, and the
static equations give global Schwarzschild--anti-de Sitter identification.
The remaining reconstruction recovers both original tensors and the proper
hypersurface embedding, including its future horizon attachment and
conformal infinity.
